\chapter{Contrato de Dados entre Agentes}
\label{apend:contrato}

A comunicação entre os agentes do \textit{pipeline} ocorre por meio de arquivos. O artefato central é o documento JSON produzido pelo Agente~2, que serve de \textbf{contrato} entre a identificação de personas (Agente~2) e a geração do documento de requisitos (Agente~3). Este apêndice apresenta a estrutura desse documento. Os campos detalhados de \texttt{personas} e de \texttt{user\_stories} seguem os esquemas reproduzidos no Apêndice~\ref{apend:prompts}.

{\footnotesize
\begin{verbatim}
{
  "metadata": {
    "arquivo_origem": "nome do arquivo de transcrição de entrada",
    "data_processamento": "carimbo de data/hora do processamento",
    "modelo_usado": "modelo de linguagem utilizado (ex.: claude-haiku-4-5)"
  },
  "analise_personas": {
    "personas": [ /* atores do sistema (ver Apêndice A) */ ],
    "participantes_reuniao": [ /* quem falou na reunião (metadado) */ ],
    "resumo_conversa": "resumo do sistema/funcionalidade discutidos",
    "tipo_interacao": "ex.: reunião de levantamento de requisitos",
    "gestao_projeto": { /* decisões de gestão (não são requisitos) */ }
  },
  "historias_usuario": {
    "user_stories": [ /* histórias por persona, no formato INVEST */ ]
  }
}
\end{verbatim}
}

A separação explícita entre \texttt{personas} (atores do sistema) e \texttt{participantes\_reuniao} (quem efetivamente falou na reunião) é uma decisão de projeto deliberada: ela evita que participantes da reunião sejam tratados, por engano, como atores do software --- um modo de falha observado em versões anteriores do sistema. De modo análogo, o campo \texttt{gestao\_projeto} isola as discussões de gestão (cronograma, prazos, alocação de equipe), de forma que elas não contaminem os requisitos funcionais derivados a jusante.
